\chapter{1\up{ères} lignes peu ou prou communes}
\begin{mintedbox}{python}
# -*- coding: utf-8 -*-
"""
Created on Mon Dec 19 19:50:23 2022

@author: adaml
"""


import numpy as np
import matplotlib.pyplot as plt
import matplotlib.animation
import matplotlib.patches
import matplotlib.transforms
import matplotlib.lines
import scipy.interpolate
import shapely
import shapely.geometry
import time
from scipy.integrate import quad as integrale


VITESSE_INIT_VOIT = 5  # vitesse initiale du véhicule
VITESSE_MAX_VOIT = 15  # vitesse maximale (pour rester dans les zones linéaires des modèles)
ACCEL_LAT_MAX = 1.5  # accélération latérale maximale (pour rester dans les zones linéaires des modèles)
LONGUEUR_VEHIC = 5  # longueur du véhicule
LARGEUR_VEHIC = 2  # largeur du véhicule
POSITION_VOITURE_INIT = (-10, -2)  # position (x, y) initiale de la voiture
INCLINAISON_VOITURE_INIT = 0  # inclinaison initiale de la voiture (en degrés)
BRAQUAGE_VOLANT_INIT = 0  # angle volant/voiture en degrés
PRECISION = 15  # nombre de points définissant l'arc de cercle. ie nombre de positions où la voiture ira
TEMPS_REAC = 90  # temps de génération de tentacules
CONE_BRAQUAGE = 0.02  # valeur de braquage définissant des trajectoires "lisses". cf get_tentacules_navigables_et_lisses
COMMENTAIRES = True  # si des commentaires sur l'évolution de l'algo doivent être affichés dans la console
PETIT_POUCET = True  # si on veut récupérer la trajectoire de la voiture
SECURITE = True  # vérifie les collisions. False si aucune sécu, None si check sans arrêter, True si arrêt lorsque collision

# courbure = ro
# courbure = 1/rayon dans le cas d'un cercle
# braquage = delta = arctan(L/rayon) quasi logique (schéma + perpendiculaire)
# longueur = angle*rayon
# empattement +-= longueur voiture


"""
SOME CODE
"""

\end{mintedbox}
